\documentclass[11pt,a4paper]{article}

% ============================================================
% Packages (core only — texlive-latex-base + texlive-base)
% ============================================================
\usepackage{amsmath, amssymb, amsthm}
\usepackage[a4paper, margin=2.5cm]{geometry}
\usepackage{array}

% ============================================================
% Theorem environments
% ============================================================
\theoremstyle{definition}
\newtheorem{definition}{Definition}[section]
\newtheorem{remark}{Remark}[section]
\newtheorem{proposition}{Proposition}[section]

\theoremstyle{plain}
\newtheorem{theorem}{Theorem}[section]
\newtheorem{lemma}{Lemma}[section]

% ============================================================
% Notation macros
% ============================================================
\newcommand{\R}{\mathbb{R}}
\newcommand{\E}{\mathbb{E}}
\newcommand{\N}{\mathbb{N}}
\newcommand{\Xmat}{\mathbf{X}}
\newcommand{\Umat}{\mathbf{U}}
\newcommand{\Smat}{\boldsymbol{\Sigma}}
\newcommand{\Vmat}{\mathbf{V}}
\newcommand{\bvec}[1]{\mathbf{#1}}
\newcommand{\norm}[1]{\left\lVert #1 \right\rVert}
\newcommand{\inner}[2]{\langle #1,\, #2 \rangle}
\newcommand{\corr}{\operatorname{corr}}
\newcommand{\med}{\operatorname{median}}
\newcommand{\diag}{\operatorname{diag}}
\newcommand{\rank}{\operatorname{rank}}
\newcommand{\coloneqq}{:=}
\DeclareMathOperator*{\argmax}{arg\,max}
\DeclareMathOperator*{\argmin}{arg\,min}

% Algorithmic environment (inline, no package needed)
\newcounter{algline}
\newenvironment{pseudocode}[1]{%
  \medskip\noindent\textbf{Algorithm:} #1\par\smallskip
  \setcounter{algline}{0}%
  \begin{list}{}{%
    \setlength{\topsep}{0pt}%
    \setlength{\itemsep}{1pt}%
    \setlength{\parsep}{0pt}%
    \setlength{\leftmargin}{2em}%
  }%
}{\end{list}\medskip}
\newcommand{\aline}[1]{\item \stepcounter{algline}\textbf{\thealgline.}\ #1}
\newcommand{\aindent}[1]{\item \hspace{1.5em}#1}

% ============================================================
\begin{document}

\title{\textbf{Event-Aligned Singular Value Decomposition\\[4pt]
for Data-Driven Pattern Kernel Discovery\\[4pt]
in Cryptocurrency Perpetual Futures}}

\author{Onat Y\i lmaz\\
\textit{Independent Researcher}\\
\texttt{yogt1984@gmail.com}}
\date{May 2026}
\maketitle

% ============================================================
% Abstract
% ============================================================

\begin{abstract}
We propose a data-driven method for discovering predictive multi-candle
price patterns in high-frequency cryptocurrency markets. The pipeline
defines six analytical event types (breakouts, turtle-soup false breakouts,
bull/bear traps), extracts OHLCV windows aligned to each event, decomposes
each window into four signal channels (body, upper wick, lower wick,
volume), and applies Singular Value Decomposition to discover the dominant
shape variations. A subsequent Information Coefficient gate with
Benjamini--Hochberg false discovery rate control identifies which discovered
shapes correlate with forward returns. The method inherits the
interpretability of classical technical analysis (events are analytically
defined) while leveraging the statistical efficiency of data-driven basis
discovery (shapes are learned, not hand-designed). Crucially, SVD operates
on price geometry alone without access to forward returns, providing a
structural safeguard against overfitting. Applying the pipeline to BTC
perpetual futures tick data (7.2 million ticks, 12{,}059 one-minute
candles), we discover 6 statistically significant pattern kernels from 158
candidates after FDR correction ($q = 0.05$). Turtle-soup and trap event
types produce the strongest signals ($|\text{IC}| \in [0.20, 0.36]$), while
breakout patterns show no predictive power after multiple testing
correction. Rolling SVD stability analysis confirms that the discovered
shapes are temporally persistent ($\bar{\rho} > 0.85$ for the strongest
kernels), supporting their use as real-time features via cosine similarity
scoring.
\end{abstract}

\smallskip
\noindent\textbf{Keywords:} pattern recognition, singular value
decomposition, market microstructure, technical analysis, cryptocurrency,
matched filtering, event-aligned averaging

\smallskip
\noindent\textbf{JEL Classification:} C14, C58, G14, G17

\tableofcontents
\bigskip

% ============================================================
\section{Introduction}
\label{sec:intro}
% ============================================================

\subsection{The Problem}

Classical technical analysis assigns fixed geometric shapes to named
patterns: a ``head and shoulders'' has a specific form; a ``hammer'' has a
long lower wick and a small body. These hand-designed templates carry two
fundamental limitations: (1) the actual shape of a pattern in a given market
and regime may differ from the textbook version, and (2) there is no
principled mechanism for deciding \emph{which} shape variation predicts
future returns as opposed to merely being visually recognisable.

This paper addresses both limitations. We define market events
\emph{analytically} by price action conditions (Section~\ref{sec:events}),
making event detection unambiguous and replicable. The shapes associated
with those events are then \emph{discovered from data} via Singular Value
Decomposition (Section~\ref{sec:svd}). Only the discovered shapes that
survive a statistical Information Coefficient gate
(Section~\ref{sec:ic}) are retained. The result is a kernel library that
is both data-driven and statistically validated.

\subsection{Analogues in Other Disciplines}

The methodology draws on two established signal processing paradigms.

\paragraph{Event-Related Potentials (ERP) in neuroscience.}
An ERP is computed by locking EEG recordings to a stimulus event and
averaging across trials. The average converges to the signal phase-locked to
the event, while trial-to-trial noise cancels. Our event-aligned matrix
$\Xmat^{(e,c)}$ plays exactly this role: each row is one ``trial'' (one
occurrence of the event), and SVD generalises simple averaging by retaining
not just the mean shape but the full low-dimensional subspace of shape
variation. The left singular vector $\bvec{u}_1$ encodes trial-to-trial
amplitude variation; the right singular vector $\bvec{v}_1$ encodes the
dominant waveform shape. This connection to neuroscience ERP methodology is,
to our knowledge, novel in the financial pattern recognition literature.

\paragraph{Matched filtering in radar and seismology.}
A matched filter maximises the signal-to-noise ratio for detecting a known
template $\bvec{k}$ in a noisy signal $\bvec{x}(t)$ by computing the
cross-correlation $s(t) = \bvec{k}^\top \bvec{x}(t)$ (Turin, 1960). The
SVD right singular vectors $\bvec{v}_k$ serve as data-driven templates, and
the online scoring step (Section~\ref{sec:online}) is exactly a matched
filter bank---one filter per surviving component.

\subsection{Key Insight: Shape Discovery Before Return Fitting}

The critical anti-overfitting property is the ordering of operations:
\begin{enumerate}
    \item SVD discovers basis functions using only the \emph{event-aligned
          price data}. It has no access to forward returns.
    \item IC filtering checks which of these already-discovered shapes
          correlate with forward returns.
\end{enumerate}
We are not fitting shapes to returns. We are asking whether naturally
occurring shape variation---variation the market itself exhibits around
analytically-defined events---is predictive. This ordering ensures that the
discovered kernels are interpretable geometric objects, not overfit
regression artefacts.

\subsection{Contributions}

\begin{enumerate}
    \item A complete pipeline for data-driven pattern kernel discovery that
          combines analytical event definitions, multi-channel OHLCV
          decomposition, SVD basis discovery, and IC-based statistical
          filtering with FDR control.
    \item A multi-channel candlestick representation (body, upper wick,
          lower wick, volume) that enables SVD to capture shape variation
          independently across price geometry and volume dynamics.
    \item Empirical evidence from BTC perpetual futures showing that
          turtle-soup and trap event types produce temporally stable,
          statistically significant pattern kernels, while breakout patterns
          do not survive multiple testing correction.
    \item An analytical basis dictionary (10 functions) for interpretability
          diagnostics and warm-start initialisation, together with evidence
          that data-driven shapes are genuinely novel (maximum alignment
          $< 0.5$ with any analytical basis).
\end{enumerate}

% ============================================================
\section{Related Work}
\label{sec:related}
% ============================================================

\paragraph{Technical pattern recognition.}
Lo, Mamaysky, and Wang (2000) provided the first rigorous statistical
evaluation of technical analysis, using kernel regression to detect
classical patterns (head-and-shoulders, triangles, double tops) and testing
their information content via bootstrap. Their approach identifies
predefined pattern types; ours discovers shapes endogenously via SVD.
Leigh \emph{et al.} (2002) used template matching with fixed bull-flag
templates. Chang and Osler (1999) studied the profitability of
head-and-shoulders patterns in foreign exchange. All of these methods rely
on hand-designed templates, which our approach replaces with data-driven
discovery.

\paragraph{PCA and random matrix theory in finance.}
The application of eigenvalue decomposition to financial correlation
matrices is well established (Laloux \emph{et al.}, 1999; Plerou \emph{et
al.}, 2002; Bouchaud and Potters, 2003). These works decompose
cross-sectional return covariance to separate signal from noise. Our
application differs fundamentally: we apply SVD to \emph{time-domain}
event-aligned windows of a single asset's candle data, decomposing temporal
shape variation rather than cross-sectional correlation structure.

\paragraph{Event studies.}
The event study methodology (Fama \emph{et al.}, 1969; MacKinlay, 1997)
examines abnormal returns around corporate events (earnings, mergers). The
standard approach averages returns across events; it does not examine the
\emph{shape} of price action preceding the event. Our event-aligned matrix
construction is analogous to cumulative abnormal return (CAR) windows but
applied to raw OHLCV channels and decomposed via SVD rather than averaged.

\paragraph{Signal processing in finance.}
Cont (2001) introduced the concept of empirical scaling laws and statistical
properties of financial time series. Gatheral and Oomen (2010) connected
microstructure noise to realised volatility estimation. The matched
filtering analogy we draw on originates in radar signal detection (Turin,
1960) and gravitational wave astronomy (Allen \emph{et al.}, 2012), where
template banks of known waveforms are correlated against noisy detector
output. Our contribution adapts this framework by constructing the template
bank from data via SVD rather than from physical models.

\paragraph{Cryptocurrency microstructure.}
Makarov and Schoar (2020) studied price formation and arbitrage in
cryptocurrency markets. Ait-Sahalia and Brunetti (2024) analysed order flow
toxicity in Bitcoin markets. Our work contributes to this literature by
providing a systematic method for extracting predictive candle-level
patterns from high-frequency cryptocurrency data.

% ============================================================
\section{Notation}
\label{sec:notation}
% ============================================================

\begin{center}
\small
\begin{tabular}{lll}
\hline
\textbf{Symbol} & \textbf{Type / Range} & \textbf{Meaning} \\
\hline
$t$               & $t \in \N$                    & Discrete candle time index \\
$O_t, H_t, L_t, C_t$ & $\R_{>0}$              & Open, High, Low, Close prices at bar $t$ \\
$V_t$             & $\R_{\geq 0}$                 & Traded volume at bar $t$ \\
$b_t, u_t, l_t, v_t$ & $\R$                    & Body, upper wick, lower wick, volume channels \\
$W$               & $W \in \N$                    & Window width in candles \\
$N$               & $N \in \N$                    & Lookback for event detection \\
$K$               & $K \in \N$                    & Confirmation lag (trap events) \\
$\alpha$          & $\alpha > 0$                  & ATR multiplier for trap threshold \\
$n_e$             & $n_e \in \N$                  & Number of events of type $e$ \\
$\mathcal{E}$     & set                           & Set of event types (6 elements) \\
$\mathcal{C}$     & $\{b, u, l, v\}$              & Set of signal channels \\
$\tau$            & $\{0,\ldots,W{-}1\}$          & Within-window time index \\
$\bar{\tau}$      & $[0,1]$                       & Normalised window coordinate \\
$\Xmat^{(e,c)}$   & $\R^{n_e \times W}$           & Event-aligned data matrix \\
$\Umat, \Smat, \Vmat$ & ---                       & SVD factors \\
$\sigma_k$        & $\R_{\geq 0}$                 & $k$-th singular value \\
$\bvec{v}_k$      & $\R^W$                        & $k$-th right singular vector (basis shape) \\
$\bvec{u}_k$      & $\R^{n_e}$                    & $k$-th left singular vector (event loadings) \\
$\mathrm{EVR}_k$  & $[0,1]$                       & Explained variance ratio of component $k$ \\
$\mathrm{SI}$     & $\R_{\geq 0}$                 & Stereotypy index $\sigma_1/\sigma_2$ \\
$\mathrm{IC}_k$   & $[-1,1]$                      & Information Coefficient for component $k$ \\
$r_i^{(H)}$       & $\R$                          & Forward log-return at horizon $H$ \\
$s^{(e,c)}_j(t)$  & $[-1,1]$                      & Cosine-similarity score \\
$S^{(e)}(t)$      & $\R$                          & Aggregated score for event type $e$ \\
$w_c$             & $w_c \geq 0$                  & Channel aggregation weight \\
$\mathrm{ATR}_t$  & $\R_{>0}$                     & Average True Range at time $t$ \\
$\phi_i$          & $[0,1] \to \R$                & Analytical basis function $i$ \\
\hline
\end{tabular}
\end{center}

% ============================================================
\section{Candlestick Decomposition into Signal Channels}
\label{sec:channels}
% ============================================================

\subsection{OHLCV Decomposition}

A single OHLCV candle at time $t$ is the tuple $(O_t, H_t, L_t, C_t, V_t)$
satisfying
\begin{equation}
    L_t \leq \min(O_t, C_t) \leq \max(O_t, C_t) \leq H_t, \quad V_t \geq 0.
    \label{eq:ohlcv_constraint}
\end{equation}
We decompose each candle into four \emph{signal channels}.

\begin{definition}[Signal Channels]
\label{def:channels}
For a candle at time $t$, define:
\begin{align}
    b_t &:= C_t - O_t
          \label{eq:body} \\
    u_t &:= H_t - \max(C_t,\, O_t)
          \label{eq:upper_wick} \\
    l_t &:= \min(C_t,\, O_t) - L_t
          \label{eq:lower_wick} \\
    v_t &:= V_t
          \label{eq:volume_channel}
\end{align}
\end{definition}

\begin{remark}
By \eqref{eq:ohlcv_constraint}, $u_t \geq 0$ and $l_t \geq 0$ always. The
body $b_t$ is signed: $b_t > 0$ for a bull candle, $b_t < 0$ for a bear
candle. The decomposition is invertible: given $(b_t, u_t, l_t, O_t)$ one
recovers $(O_t, H_t, L_t, C_t)$ exactly. The channel set is
$\mathcal{C} = \{b, u, l, v\}$.
\end{remark}

The multi-channel representation is essential: different event types exhibit
their characteristic signatures in different channels. As we show
empirically (Section~\ref{sec:results}), the volume channel carries
independent predictive information for trap events that is absent from the
price channels alone.

\subsection{Window Extraction}

\begin{definition}[Event Window]
\label{def:window}
For an event detected at time $t_e$ and window width $W \in \N$, the
\emph{event window} of channel $c \in \mathcal{C}$ is the vector
\begin{equation}
    \bvec{x}^{(c)}(t_e) :=
    \bigl(x^{(c)}_{t_e - W + 1},\;
          x^{(c)}_{t_e - W + 2},\;
          \ldots,\;
          x^{(c)}_{t_e}\bigr)^\top
    \in \R^W,
    \label{eq:window_vector}
\end{equation}
where $x^{(c)}_t \in \{b_t, u_t, l_t, v_t\}$ for $c \in \{b, u, l, v\}$
respectively.
\end{definition}

\subsection{Per-Window Normalisation}
\label{sec:normalisation}

Raw price levels are not comparable across events separated in time or
across volatility regimes. We apply a \emph{shape-preserving} normalisation
that removes level and scale but retains geometry.

\begin{definition}[ATR Normalisation]
\label{def:normalisation}
Let $\mathrm{ATR}_t$ be the $N_{\mathrm{atr}}$-period Average True Range
(Wilder, 1978):
\begin{equation}
    \mathrm{ATR}_t :=
    \frac{1}{N_{\mathrm{atr}}}
    \sum_{j=0}^{N_{\mathrm{atr}}-1}
    \max\!\bigl(H_{t-j} - L_{t-j},\;
                |H_{t-j} - C_{t-j-1}|,\;
                |L_{t-j} - C_{t-j-1}|\bigr).
    \label{eq:atr}
\end{equation}
The normalised window for event $i$ with event time $t_e^{(i)}$ is
\begin{equation}
    \tilde{x}^{(c)}_i(\tau) :=
    \frac{x^{(c)}_{t_e^{(i)} - W + 1 + \tau} - \mu^{(c)}_i}
         {\mathrm{ATR}_{t_e^{(i)}}},
    \quad \tau = 0, 1, \ldots, W-1,
    \label{eq:normalised_window}
\end{equation}
where $\mu^{(c)}_i := \frac{1}{W}\sum_{\tau=0}^{W-1}
x^{(c)}_{t_e^{(i)}-W+1+\tau}$ is the within-window mean.
\end{definition}

% ============================================================
\section{Analytical Event Definitions}
\label{sec:events}
% ============================================================

We define six event types in $\mathcal{E}$. Let $N$ be the lookback period
and let $\mathrm{ATR}_t$ be as in \eqref{eq:atr}.

\subsection{Breakout Events}

\begin{definition}[Breakout, Bull]
\label{def:breakout_bull}
A \emph{bull breakout} occurs at time $t$ if and only if
\begin{equation}
    C_t > \max_{j=1}^{N} H_{t-j}
    \quad \text{and} \quad
    V_t > 1.5 \cdot \med\!\bigl(V_{t-N}, \ldots, V_{t-1}\bigr).
    \label{eq:breakout_bull}
\end{equation}
\end{definition}

\begin{definition}[Breakout, Bear]
\label{def:breakout_bear}
A \emph{bear breakout} occurs at time $t$ if and only if
\begin{equation}
    C_t < \min_{j=1}^{N} L_{t-j}
    \quad \text{and} \quad
    V_t > 1.5 \cdot \med\!\bigl(V_{t-N}, \ldots, V_{t-1}\bigr).
    \label{eq:breakout_bear}
\end{equation}
\end{definition}

The volume condition filters genuine momentum moves from slow drift: a real
breakout typically registers at least $1.5\times$ the recent median volume.
The $N$-period range maximum $\max_{j=1}^{N} H_{t-j}$ is the upper
Donchian channel (Donchian, 1960), excluding the current bar.

\subsection{Turtle Soup Events (False Range Break)}

\emph{Turtle Soup} refers to a pattern described by Raschke and Connors
(1995): price transiently penetrates an $N$-period extreme, then closes
back inside the prior range, suggesting that the breakout attracted
stop-loss orders absorbed by a larger counterparty.

\begin{definition}[Turtle Soup, Bull]
\label{def:ts_bull}
A \emph{bull turtle soup} (fake breakdown, anticipates reversal long)
occurs at time $t$ if and only if
\begin{equation}
    L_t < \min_{j=1}^{N} L_{t-j}
    \quad \text{and} \quad
    C_t > \min_{j=1}^{N} L_{t-j}.
    \label{eq:ts_bull}
\end{equation}
\end{definition}

\begin{definition}[Turtle Soup, Bear]
\label{def:ts_bear}
A \emph{bear turtle soup} (fake breakout, anticipates reversal short)
occurs at time $t$ if and only if
\begin{equation}
    H_t > \max_{j=1}^{N} H_{t-j}
    \quad \text{and} \quad
    C_t < \max_{j=1}^{N} H_{t-j}.
    \label{eq:ts_bear}
\end{equation}
\end{definition}

In both cases the candle's extreme pierces the historical range but the
close reverts. The lower wick channel $l_t$ is expected to be especially
large for Definition~\ref{def:ts_bull}, providing a signal-rich feature for
SVD.

\subsection{Trap Events (Confirmed False Breakout)}

A trap event requires multi-bar confirmation: the breakout occurs at time
$t$, and the close at time $t + K$ must have reversed by at least
$\alpha \cdot \mathrm{ATR}_t$.

\begin{definition}[Bull Trap]
\label{def:bull_trap}
A \emph{bull trap} is a pair $(t, t+K)$ satisfying:
\begin{itemize}
    \item[(i)] Condition \eqref{eq:breakout_bull} holds at time $t$.
    \item[(ii)] $C_{t+K} < C_t - \alpha \cdot \mathrm{ATR}_t$.
\end{itemize}
The event window is centred on the breakout candle at time $t$.
\end{definition}

\begin{definition}[Bear Trap]
\label{def:bear_trap}
A \emph{bear trap} is a pair $(t, t+K)$ satisfying:
\begin{itemize}
    \item[(i)] Condition \eqref{eq:breakout_bear} holds at time $t$.
    \item[(ii)] $C_{t+K} > C_t + \alpha \cdot \mathrm{ATR}_t$.
\end{itemize}
\end{definition}

\begin{remark}[Look-ahead bias]
\label{rem:lookahead}
Trap events use information from $t+K$ to label the event at $t$. They are
therefore only valid for offline kernel discovery. The aligned window is
extracted at $t$, and the confirmation at $t+K$ partitions events into
``trap'' versus ``genuine breakout'' classes. Deploying the kernel in real
time introduces no look-ahead, because the kernel scores the current
window at bar $t$ without knowledge of future prices.
\end{remark}

\paragraph{Parameter defaults.}
\begin{center}
\begin{tabular}{llll}
\hline
\textbf{Parameter} & \textbf{Symbol} & \textbf{Default} & \textbf{Valid Range} \\
\hline
Range lookback           & $N$                   & 20 bars   & $[10, 60]$ \\
Volume multiplier        & ---                   & 1.5       & $[1.2, 2.5]$ \\
ATR period               & $N_{\mathrm{atr}}$    & 14 bars   & $[7, 28]$ \\
Trap confirmation lag    & $K$                   & 3 bars    & $[1, 10]$ \\
ATR reversal threshold   & $\alpha$              & 1.0       & $[0.5, 2.0]$ \\
\hline
\end{tabular}
\end{center}

% ============================================================
\section{Event-Aligned Matrix Construction}
\label{sec:matrix}
% ============================================================

\begin{definition}[Event-Aligned Matrix]
\label{def:event_matrix}
Let $\mathcal{T}_e = \{t_e^{(1)}, t_e^{(2)}, \ldots, t_e^{(n_e)}\}$ be
the ordered set of detection times for event type $e \in \mathcal{E}$. For
channel $c \in \mathcal{C}$, the \emph{event-aligned matrix} is
\begin{equation}
    \Xmat^{(e,c)} :=
    \begin{pmatrix}
        \tilde{x}^{(c)}_1(0) & \tilde{x}^{(c)}_1(1) & \cdots & \tilde{x}^{(c)}_1(W-1) \\
        \tilde{x}^{(c)}_2(0) & \tilde{x}^{(c)}_2(1) & \cdots & \tilde{x}^{(c)}_2(W-1) \\
        \vdots & \vdots & \ddots & \vdots \\
        \tilde{x}^{(c)}_{n_e}(0) & \tilde{x}^{(c)}_{n_e}(1) & \cdots &
        \tilde{x}^{(c)}_{n_e}(W-1)
    \end{pmatrix}
    \in \R^{n_e \times W},
    \label{eq:event_matrix}
\end{equation}
where $\tilde{x}^{(c)}_i(\tau)$ is given by \eqref{eq:normalised_window}.
\end{definition}

Each row is one normalised event window. Each column corresponds to a
relative time offset $\tau$ within the window. We require
\begin{equation}
    n_e \geq n_{\min} = 100
    \label{eq:min_events}
\end{equation}
before proceeding to decomposition. Under a spiked covariance model, the
variance of the estimated right singular vector is
$O(\sigma^2 W / (n_e \sigma_k^2))$; for $W = 20$ and typical signal-to-noise
ratios, $n_e = 100$ achieves $O(10\%)$ estimation error.

% ============================================================
\section{SVD Decomposition: Basis Discovery}
\label{sec:svd}
% ============================================================

\subsection{The Decomposition}

\begin{definition}[Thin SVD]
\label{def:svd}
For $\Xmat := \Xmat^{(e,c)} \in \R^{n_e \times W}$ with $n_e \geq W$,
the thin Singular Value Decomposition is
\begin{equation}
    \Xmat = \Umat \, \Smat \, \Vmat^\top,
    \label{eq:svd}
\end{equation}
where
$\Umat \in \R^{n_e \times W}$ with $\Umat^\top \Umat = \mathbf{I}_W$,
$\Smat = \diag(\sigma_1, \ldots, \sigma_W)$ with
$\sigma_1 \geq \cdots \geq \sigma_W \geq 0$, and
$\Vmat \in \R^{W \times W}$ with $\Vmat^\top \Vmat = \mathbf{I}_W$.
\end{definition}

The columns $\bvec{v}_k = \Vmat_{:,k} \in \R^W$ are the \emph{discovered
basis functions} (shapes), ordered from dominant to residual. The columns
$\bvec{u}_k = \Umat_{:,k} \in \R^{n_e}$ give the \emph{per-event loading}
on shape $k$: the scalar $u_{ik}$ measures how strongly event $i$ projects
onto shape $k$.

\subsection{Interpretive Quantities}

\begin{definition}[Explained Variance Ratio]
\label{def:evr}
\begin{equation}
    \mathrm{EVR}_k :=
    \frac{\sigma_k^2}{\displaystyle\sum_{j=1}^{W} \sigma_j^2}
    = \frac{\sigma_k^2}{\norm{\Xmat}_F^2}.
    \label{eq:evr}
\end{equation}
\end{definition}

\begin{definition}[Stereotypy Index]
\label{def:stereotypy}
The \emph{stereotypy index} of event class $(e, c)$ is
\begin{equation}
    \mathrm{SI}^{(e,c)} := \frac{\sigma_1}{\sigma_2}.
    \label{eq:stereotypy}
\end{equation}
A high stereotypy index ($\mathrm{SI} \gg 1$) indicates a single dominant
shape; an index near 1 indicates a broad, multi-modal shape distribution.
\end{definition}

\subsection{Truncation Rule}

We retain the top $K^{(e,c)}$ components where cumulative explained
variance first exceeds $\rho^* = 0.80$:
\begin{equation}
    K^{(e,c)} :=
    \min\!\left\{
        K \in \{1, \ldots, W\} \,:\,
        \sum_{k=1}^{K} \mathrm{EVR}_k \geq \rho^*
    \right\}.
    \label{eq:truncation}
\end{equation}

By the Eckart--Young--Mirsky theorem (Eckart and Young, 1936), the retained
components provide the best rank-$K$ approximation to $\Xmat$ in Frobenius
norm, with relative reconstruction error
$\epsilon_K = \sqrt{1 - \sum_{k=1}^{K} \mathrm{EVR}_k}$.

% ============================================================
\section{Predictive Filtering: The IC Gate}
\label{sec:ic}
% ============================================================

\subsection{Forward Returns}

\begin{definition}[Forward Log-Return]
\label{def:forward_return}
For event $i$ at time $t_e^{(i)}$ and horizon $H$ candles:
\begin{equation}
    r_i^{(H)} := \ln\!\frac{C_{t_e^{(i)}+H}}{C_{t_e^{(i)}}}.
    \label{eq:forward_return}
\end{equation}
The return begins strictly after the event window ends, ensuring no overlap
between the information used by SVD and the return being predicted.
\end{definition}

\subsection{Information Coefficient}

\begin{definition}[Information Coefficient, IC]
\label{def:ic}
For SVD component $k$ of event type $(e, c)$:
\begin{equation}
    \mathrm{IC}_k^{(e,c,H)}
    := \corr\!\bigl(\bvec{u}_k,\; \bvec{r}^{(H)}\bigr).
    \label{eq:ic}
\end{equation}
\end{definition}

Component $k$ is retained if $|\mathrm{IC}_k^{(e,c,H)}| \geq
\delta_{\mathrm{IC}} = 0.03$.

\begin{remark}[Why this is not overfitting]
The loadings $\bvec{u}_k$ are determined entirely by the price-action
matrix $\Xmat^{(e,c)}$. The SVD objective minimises reconstruction error
with no knowledge of $\bvec{r}^{(H)}$. The IC computation is a
\emph{test} on the SVD output, not a training step. This is analogous to
computing the correlation between a PCA projection and an external label.
\end{remark}

\subsection{Multiple Comparison Correction}

Across all combinations $(e, c, k, H)$, the IC tests constitute a family
of hypotheses. We apply Benjamini--Hochberg FDR control (Benjamini and
Hochberg, 1995) at level $q = 0.05$:
\begin{enumerate}
    \item Compute the two-sided $p$-value under $H_0\colon \mathrm{IC} = 0$
          using the $t$-statistic
          \begin{equation}
              t_k = \mathrm{IC}_k
                    \sqrt{\frac{n_e - 2}{1 - \mathrm{IC}_k^2}},
              \quad t_k \sim t_{n_e-2} \text{ under } H_0.
              \label{eq:ic_tstat}
          \end{equation}
    \item Order $p$-values: $p_{(1)} \leq p_{(2)} \leq \cdots \leq p_{(m)}$.
    \item Reject $H_0$ for all $j \leq j^*$, where
          $j^* := \max\{j : p_{(j)} \leq jq/m\}$.
\end{enumerate}

% ============================================================
\section{Kernel Construction and Online Scoring}
\label{sec:online}
% ============================================================

\subsection{Kernel Library}

After IC filtering, the kernel library is
\begin{equation}
    \mathcal{K} :=
    \bigl\{\bvec{k}^{(e,c)}_j \in \R^W \,:\,
    \text{$(e,c,j)$ survives IC gate and BH-FDR}\bigr\},
    \label{eq:kernel_library}
\end{equation}
where $\bvec{k}^{(e,c)}_j := \bvec{v}^{(e,c)}_j$. Each kernel is unit
norm, orthogonal to other kernels within the same $(e,c)$ class, and
mean-free by construction.

\subsection{Cosine Similarity Scoring}

At production time $t$, the current normalised window for channel $c$ is
$\bvec{x}^{(c)}_t \in \R^W$. The cosine similarity score is
\begin{equation}
    s^{(e,c)}_j(t) :=
    \frac{\inner{\bvec{x}^{(c)}_t}{\bvec{k}^{(e,c)}_j}}
         {\norm{\bvec{x}^{(c)}_t}_2},
    \label{eq:cosine_score}
\end{equation}
where the simplification uses $\norm{\bvec{k}}_2 = 1$. The score satisfies
$s \in [-1, 1]$ by the Cauchy--Schwarz inequality.

\subsection{Multi-Channel Aggregation}

For event type $e$, the composite score aggregates across channels and
surviving components:
\begin{equation}
    S^{(e)}(t) :=
    \sum_c w_c \sum_{\{j : \bvec{k}^{(e,c)}_j \in \mathcal{K}\}}
    \mathrm{IC}_j \cdot s^{(e,c)}_j(t),
    \label{eq:composite_score}
\end{equation}
where $w_c$ are channel weights ($w_c = 0.25$ by default or
$w_c \propto |\mathrm{IC}|$ for IC-weighted aggregation). The IC weight
ensures that components with higher predictive power contribute more.

The inner product $\inner{\bvec{x}}{\bvec{k}}$ requires $O(W)$ operations
per kernel per bar. For $|\mathcal{K}| \leq 20$ and $W = 20$, online
scoring costs 400 multiply-accumulate operations per candle---well within
real-time constraints.

% ============================================================
\section{Validation Framework}
\label{sec:validation}
% ============================================================

\subsection{Walk-Forward Validation}

\begin{pseudocode}{Walk-Forward Kernel Discovery and Validation}
\aline{Split dataset at $T_{\mathrm{train}} = \lfloor 0.60 \cdot T \rfloor$.}
\aline{\textbf{Discovery} on $t \leq T_{\mathrm{train}}$:}
\aindent{For each $(e, c)$: detect events, build $\Xmat^{(e,c)}$, SVD, IC gate.}
\aline{Apply BH-FDR across all in-sample tests.}
\aline{\textbf{Validation} on $t > T_{\mathrm{train}}$:}
\aindent{For each OOS event: score with $\mathcal{K}$; record $S^{(e)}$ and $r^{(H)}$.}
\aline{Compute OOS IC; report IC decay ratio.}
\end{pseudocode}

A kernel is considered robust if
\begin{equation}
    \rho_{\mathrm{decay}} :=
    \frac{\mathrm{IC}^{\mathrm{OOS}}_k}{\mathrm{IC}^{\mathrm{IS}}_k}
    \geq 0.5.
    \label{eq:ic_decay}
\end{equation}

\subsection{Rolling SVD Stability}

\begin{definition}[Rolling Stability]
\label{def:rolling_stability}
For rolling windows $\mathcal{D}_1, \ldots, \mathcal{D}_M$ with overlap,
let $\bvec{v}^{(m)}_1$ be the first right singular vector on window $m$.
The stability between consecutive windows is
\begin{equation}
    \rho_m :=
    \bigl|\inner{\bvec{v}^{(m)}_1}{\bvec{v}^{(m+1)}_1}\bigr|.
    \label{eq:rolling_stability}
\end{equation}
The overall stability is $\bar{\rho} = (M-1)^{-1}\sum_{m=1}^{M-1} \rho_m$.
We require $\bar{\rho} \geq 0.70$.
\end{definition}

% ============================================================
\section{Analytical Basis Dictionary}
\label{sec:basis}
% ============================================================

As a diagnostic, we define an analytical basis dictionary against which
data-driven shapes can be compared. Let
$\bar{\tau} = \tau/(W-1) \in [0,1]$.

\paragraph{Polynomial family.}
\begin{align}
    \phi_{q1}(\bar{\tau}) &= \bar{\tau}^2, \quad
    \phi_{q2}(\bar{\tau}) = (\bar{\tau} - \tfrac{1}{2})^2, \quad
    \phi_{q3}(\bar{\tau}) = (1 - \bar{\tau})^2.
    \label{eq:poly_basis}
\end{align}

\paragraph{Cubic family.}
\begin{align}
    \phi_{c1}(\bar{\tau}) &= \bar{\tau}^3, \quad
    \phi_{c2}(\bar{\tau}) = (\bar{\tau} - \tfrac{1}{2})^3, \quad
    \phi_{c3}(\bar{\tau}) = (1 - \bar{\tau})^3, \quad
    \phi_{s1}(\bar{\tau}) = 3\bar{\tau}^2 - 2\bar{\tau}^3.
    \label{eq:cubic_basis}
\end{align}
The smoothstep $\phi_{s1}$ (Hermite easing function) satisfies
$\phi_{s1}(0)=0$, $\phi_{s1}(1)=1$, $\phi'_{s1}(0)=\phi'_{s1}(1)=0$.

\paragraph{Exponential family} ($a \in [3,8]$).
\begin{align}
    \phi_{e1}(\bar{\tau}) &= e^{-a\bar{\tau}}, \quad
    \phi_{e2}(\bar{\tau}) = e^{-a(1-\bar{\tau})}, \quad
    \phi_{b1}(\bar{\tau}) = e^{-a(\bar{\tau}-1/2)^2}.
    \label{eq:exp_basis}
\end{align}

\paragraph{Alignment diagnostic.}
For each SVD-discovered vector $\bvec{v}_k$ and analytical basis $\phi_i$
(discretised to $W$ points and $L^2$-normalised):
\begin{equation}
    \rho_{k,i} :=
    \bigl|\inner{\bvec{v}_k}{\hat{\boldsymbol{\phi}}_i}\bigr|.
    \label{eq:alignment}
\end{equation}
Values $\rho > 0.85$ indicate strong alignment (interpretable shape);
$\rho < 0.5$ indicates a genuinely novel data-driven shape.

% ============================================================
\section{Empirical Results}
\label{sec:results}
% ============================================================

We apply the pipeline to BTC perpetual futures on Hyperliquid, a
decentralised derivatives exchange, using 100\,ms tick data aggregated to
60-second candles.

\subsection{Data}

\begin{center}
\begin{tabular}{ll}
\hline
\textbf{Property} & \textbf{Value} \\
\hline
Asset              & BTC-USD perpetual \\
Exchange           & Hyperliquid \\
Tick frequency     & 100\,ms \\
Total ticks        & 7,235,623 \\
Candle frequency   & 60\,s \\
Total candles      & 12,059 \\
Parameters         & $W=20$, $N=20$, $K=3$, $\alpha=1.0$, $N_\text{atr}=14$ \\
\hline
\end{tabular}
\end{center}

\subsection{Event Detection}

\begin{center}
\begin{tabular}{lrr}
\hline
\textbf{Event Type} & \textbf{Count} & \textbf{Rate (per 1000 candles)} \\
\hline
Breakout bull  & 418 & 34.7 \\
Breakout bear  & 522 & 43.3 \\
Turtle bull    & 286 & 23.7 \\
Turtle bear    & 241 & 20.0 \\
Trap bull      & 136 & 11.3 \\
Trap bear      & 172 & 14.3 \\
\hline
\textbf{Total} & \textbf{1,775} & \textbf{147.2} \\
\hline
\end{tabular}
\end{center}

Breakout events are most frequent (bear $>$ bull asymmetry). Trap events
are rarest, consistent with their stricter multi-bar confirmation criteria.
All event types exceed the $n_{\min} = 100$ threshold required for stable
SVD.

\subsection{SVD Stereotypy}

Table~\ref{tab:stereotypy} reports the stereotypy index $\mathrm{SI} =
\sigma_1/\sigma_2$ across all (event, channel) pairs.

\begin{table}[h]
\begin{center}
\begin{tabular}{lrrrr}
\hline
\textbf{Event Type} & \textbf{body} & \textbf{upper wick} & \textbf{lower wick} & \textbf{volume} \\
\hline
Breakout bull  & 1.24 & 1.54 & 1.80 & \textbf{2.42} \\
Breakout bear  & 1.39 & 1.55 & \textbf{1.93} & 1.54 \\
Turtle bull    & 1.02 & 1.28 & \textbf{1.59} & 1.49 \\
Turtle bear    & 1.16 & 1.32 & 1.36 & \textbf{1.52} \\
Trap bull      & \textbf{1.69} & 1.55 & 1.54 & \textbf{2.74} \\
Trap bear      & 1.17 & 1.34 & 1.59 & \textbf{1.67} \\
\hline
\end{tabular}
\end{center}
\caption{Stereotypy index by event type and channel. Bold indicates highest
per row. Volume consistently shows the highest stereotypy for breakout and
trap events, indicating a more consistent volume signature than price shape.
Trap bull volume ($\mathrm{SI} = 2.74$) is the highest across all pairs.}
\label{tab:stereotypy}
\end{table}

\subsection{Surviving Kernels}

From 158 candidates passing $|\mathrm{IC}| > 0.03$, BH-FDR correction at
$q = 0.05$ retains \textbf{6 kernels} (Table~\ref{tab:survivors}).

\begin{table}[h]
\begin{center}
\small
\begin{tabular}{clllrrr}
\hline
\# & \textbf{Event} & \textbf{Channel} & \textbf{$k$} & \textbf{IC} & \textbf{$p$-value} & \textbf{EVR} \\
\hline
1 & Turtle bull  & body        & 0 & $+0.203$ & $5.4 \times 10^{-4}$ & 12.2\% \\
2 & Turtle bull  & body        & 6 & $-0.201$ & $6.3 \times 10^{-4}$ & 5.6\%  \\
3 & Turtle bear  & lower wick  & 0 & $-0.262$ & $3.7 \times 10^{-5}$ & 11.2\% \\
4 & Trap bull    & upper wick  & 6 & $+0.357$ & $2.0 \times 10^{-5}$ & 4.7\%  \\
5 & Trap bull    & volume      & 0 & $-0.270$ & $1.5 \times 10^{-3}$ & 26.5\% \\
6 & Trap bear    & body        & 7 & $-0.272$ & $3.0 \times 10^{-4}$ & 4.9\%  \\
\hline
\end{tabular}
\end{center}
\caption{Surviving kernels after BH-FDR correction. No breakout kernels
survive. Turtle and trap events dominate. The trap bull volume kernel
($k=0$, $\mathrm{EVR} = 26.5\%$) captures a quarter of all trap bull
volume variance in a single shape, confirming that volume carries
independent predictive information.}
\label{tab:survivors}
\end{table}

Three findings merit emphasis:

\begin{enumerate}
\item \textbf{No breakout kernels survive.} Despite high EVR (breakout
      shapes exist and are consistent), they carry no predictive power for
      forward returns after FDR correction. The market may have already
      incorporated breakout information by the time the pattern completes.

\item \textbf{Multi-channel contribution.} The 6 survivors span three
      channels (body, lower wick, upper wick, volume), confirming that the
      multi-channel decomposition captures information absent from the body
      channel alone.

\item \textbf{Moderate IC magnitudes.} All ICs are in the range
      $[0.20, 0.36]$, appropriate for 10-minute forward returns. The signals
      are genuine but not strong enough to trade in isolation.
\end{enumerate}

\subsection{Alignment with Analytical Basis}

All 6 surviving kernels have maximum cosine similarity $< 0.5$ against all
10 analytical basis functions. The data-driven shapes are genuinely
novel---they cannot be approximated by polynomial, cubic, or exponential
templates. This validates the SVD-based discovery approach over hand-designed
templates.

\subsection{Walk-Forward Validation}

Train: 7,235 candles (60\%); test: 4,824 candles (40\%).

\begin{center}
\small
\begin{tabular}{lllrrrr}
\hline
\textbf{Event} & \textbf{Channel} & $k$ & \textbf{IS IC} & \textbf{OOS IC} & \textbf{Decay} & \textbf{Robust} \\
\hline
Breakout bull & volume      & 0  & $+0.243$ & $-0.010$ & $-0.04$ & No \\
Breakout bear & body        & 10 & $+0.271$ & $+0.038$ & $0.14$  & No \\
Turtle bull   & body        & 7  & $+0.285$ & $+0.108$ & $0.38$  & No \\
Turtle bear   & lower wick  & 0  & $-0.420$ & $-0.108$ & $0.26$  & No \\
Trap bull     & upper wick  & 4  & $-0.474$ & $+0.126$ & $-0.27$ & No \\
Trap bull     & volume      & 0  & $-0.449$ & $+0.145$ & $-0.32$ & No \\
Trap bear     & body        & 4  & $+0.418$ & $-0.036$ & $-0.09$ & No \\
Trap bear     & body        & 5  & $+0.415$ & $+0.010$ & $0.03$  & No \\
\hline
\end{tabular}
\end{center}

No kernels meet the robustness threshold ($\rho_{\mathrm{decay}} \geq
0.50$). However, this is primarily a data volume limitation: with 12{,}059
total candles split 60/40, trap events (136--172 occurrences) yield only
$\sim$50--70 test-set events---insufficient for statistical significance.
Turtle-type kernels show partial generalisation (decay ratios 0.26--0.38
with consistent sign), suggesting that with 3--6$\times$ more data, OOS
robustness may emerge.

\subsection{Rolling Stability}

\begin{center}
\begin{tabular}{lrrl}
\hline
\textbf{Event Type} & \textbf{Mean $\bar{\rho}$} & \textbf{Min $\bar{\rho}$} & \textbf{Status} \\
\hline
Trap bull      & \textbf{0.941} & 0.865 & Stable     \\
Turtle bull    & \textbf{0.920} & 0.825 & Stable     \\
Breakout bear  & \textbf{0.855} & 0.726 & Stable     \\
Turtle bear    & 0.719          & 0.352 & Borderline \\
Breakout bull  & 0.692          & 0.165 & Unstable   \\
Trap bear      & 0.669          & 0.465 & Unstable   \\
\hline
\end{tabular}
\end{center}

The event types with the strongest IC (trap bull, turtle bull) also exhibit
the highest temporal stability ($\bar{\rho} > 0.90$). This correlation
between predictive power and shape persistence is expected: if a shape is
not stable across time, its predictive association with returns cannot be
stable either.

% ============================================================
\section{Discussion}
\label{sec:discussion}
% ============================================================

\subsection{Methodological Implications}

The central finding is that SVD applied to event-aligned candlestick
windows discovers statistically significant pattern shapes, but only for
reversal-type events (turtle soup and traps), not for momentum-type events
(breakouts). This has a natural economic interpretation: breakout patterns
are widely traded and rapidly arbitraged, so the shape of a breakout
carries no residual predictive information by the time it completes.
Turtle-soup and trap patterns, by contrast, involve deceptive price action
(false breakouts followed by reversals) that may be harder for the market
to price efficiently.

\subsection{Multi-Channel Decomposition}

The volume channel contributes the single strongest kernel (trap bull
volume, $\mathrm{EVR} = 26.5\%$, $\mathrm{IC} = -0.270$). This validates
the multi-channel approach: had we used only the body channel---as most
pattern recognition studies do---this signal would have been missed entirely.
The volume signature of a trap event is more stereotyped ($\mathrm{SI} =
2.74$) than its price signature ($\mathrm{SI} = 1.69$), suggesting that
volume dynamics carry structural information about the nature of the
breakout that is absent from price alone.

\subsection{Anti-Overfitting Architecture}

The pipeline provides three independent layers of protection:
\begin{enumerate}
\item SVD is return-blind: shapes are discovered without access to forward
      returns.
\item BH-FDR controls the expected fraction of false discoveries at
      $q = 5\%$ across all tested $(e,c,k,H)$ combinations.
\item Walk-forward OOS validation tests generalisability to unseen data.
\end{enumerate}

\subsection{Limitations}

\begin{enumerate}
\item \textbf{Data volume.} With 12{,}059 candles, rare events
      (traps: $\sim$150 occurrences) have insufficient test-set counts for
      walk-forward validation. We estimate 50{,}000+ candles
      (3--6 months at 60\,s resolution) are needed for robust OOS results.
\item \textbf{Single asset.} Results are for BTC perpetual only.
      Cross-asset universality (ETH, SOL) is untested.
\item \textbf{Single exchange.} Hyperliquid is a decentralised exchange
      with specific microstructure properties. Generalisation to centralised
      venues remains to be demonstrated.
\item \textbf{Stationarity.} The rolling stability analysis shows temporal
      persistence over the available data, but cannot guarantee persistence
      across major regime changes.
\end{enumerate}

% ============================================================
\section{Conclusion}
\label{sec:conclusion}
% ============================================================

We have presented a complete pipeline for data-driven pattern kernel
discovery in financial markets, combining analytical event definitions with
SVD-based shape decomposition and IC-gated statistical filtering. The
methodology bridges classical technical analysis and modern statistical
learning: events are defined by practitioner-interpretable price conditions,
but the shapes are discovered by the data rather than prescribed by
tradition.

Applied to BTC perpetual futures, the pipeline discovers 6 statistically
significant kernels from 158 candidates. The surviving kernels are
concentrated in turtle-soup and trap event types, with contributions from
multiple candle channels (body, wick, volume). Rolling stability analysis
confirms that the strongest kernels persist across different time windows
($\bar{\rho} > 0.90$). The discovered shapes have low alignment with
hand-designed analytical basis functions ($\max \rho < 0.5$), demonstrating
that the data-driven approach captures structures that an analyst would not
have designed a priori.

The primary limitation is data volume: walk-forward OOS validation does not
yet confirm robustness at the $\rho_{\mathrm{decay}} \geq 0.50$ threshold.
This is a data limitation rather than a methodological failure---in-sample
ICs are strong ($|\mathrm{IC}|$ up to 0.47) and turtle-type kernels show
partial OOS generalisation. Future work will extend the analysis to
multiple assets and longer time horizons.

% ============================================================
\section*{References}
% ============================================================

\begin{enumerate}
    \item Ait-Sahalia, Y.\ \& Brunetti, C.\ (2024). What do we know about
          cryptocurrency? \textit{Annual Review of Financial Economics}, 16.

    \item Allen, B.\ \emph{et al.}\ (2012). FINDCHIRP: An algorithm for
          detection of gravitational waves from inspiraling compact binaries.
          \textit{Physical Review D}, 85(12), 122006.

    \item Bandt, C.\ \& Pompe, B.\ (2002). Permutation entropy: a natural
          complexity measure for time series. \textit{Physical Review
          Letters}, 88(17), 174102.

    \item Benjamini, Y.\ \& Hochberg, Y.\ (1995). Controlling the false
          discovery rate: a practical and powerful approach to multiple
          testing. \textit{Journal of the Royal Statistical Society,
          Series B}, 57(1), 289--300.

    \item Bouchaud, J.-P.\ \& Potters, M.\ (2003). \textit{Theory of
          Financial Risk and Derivative Pricing: From Statistical Physics
          to Risk Management}, 2nd ed. Cambridge University Press.

    \item Chang, P.H.K.\ \& Osler, C.L.\ (1999). Methodical madness:
          Technical analysis and the irrationality of exchange-rate
          forecasts. \textit{Economic Journal}, 109(458), 636--661.

    \item Cont, R.\ (2001). Empirical properties of asset returns:
          stylized facts and statistical issues. \textit{Quantitative
          Finance}, 1(2), 223--236.

    \item Donchian, R.D.\ (1960). Trend-following methods in commodity
          price analysis. \textit{Commodity Yearbook.}

    \item Eckart, C.\ \& Young, G.\ (1936). The approximation of one
          matrix by another of lower rank. \textit{Psychometrika}, 1(3),
          211--218.

    \item Fama, E.F., Fisher, L., Jensen, M.C.\ \& Roll, R.\ (1969). The
          adjustment of stock prices to new information. \textit{International
          Economic Review}, 10(1), 1--21.

    \item Gatheral, J.\ \& Oomen, R.\ (2010). Zero-intelligence realized
          variance estimation. \textit{Finance and Stochastics}, 14,
          249--283.

    \item Golub, G.H.\ \& Van Loan, C.F.\ (2013). \textit{Matrix
          Computations}, 4th ed. Johns Hopkins University Press.

    \item Jolliffe, I.T.\ (2002). \textit{Principal Component Analysis},
          2nd ed.\ Springer.

    \item Laloux, L., Cizeau, P., Bouchaud, J.-P.\ \& Potters, M.\
          (1999). Noise dressing of financial correlation matrices.
          \textit{Physical Review Letters}, 83(7), 1467--1470.

    \item Leigh, W., Purvis, R.\ \& Ragusa, J.M.\ (2002). Forecasting the
          NYSE composite index with technical analysis, pattern recognizer,
          neural network, and genetic algorithm. \textit{Journal of Banking
          \& Finance}, 26(11), 2137--2154.

    \item Lo, A.W., Mamaysky, H.\ \& Wang, J.\ (2000). Foundations of
          technical analysis: computational algorithms, statistical
          inference, and empirical implementation. \textit{Journal of
          Finance}, 55(4), 1705--1765.

    \item MacKinlay, A.C.\ (1997). Event studies in economics and finance.
          \textit{Journal of Economic Literature}, 35(1), 13--39.

    \item Makarov, I.\ \& Schoar, A.\ (2020). Trading and arbitrage in
          cryptocurrency markets. \textit{Journal of Financial Economics},
          135(2), 293--319.

    \item Plerou, V., Gopikrishnan, P., Rosenow, B., Amaral, L.A.N.,
          Guhr, T.\ \& Stanley, H.E.\ (2002). Random matrix approach to
          cross correlations in financial data. \textit{Physical Review E},
          65(6), 066126.

    \item Raschke, L.B.\ \& Connors, L.A.\ (1995). \textit{Street Smarts:
          High Probability Short-Term Trading Strategies.} M.\ Gordon
          Publishing.

    \item Turin, G.L.\ (1960). An introduction to matched filters.
          \textit{IRE Transactions on Information Theory}, 6(3), 311--329.

    \item Wilder, J.W.\ (1978). \textit{New Concepts in Technical Trading
          Systems.} Trend Research.
\end{enumerate}

\end{document}
